\chapter{УДИРТГАЛ}

Сүүлийн жилүүдэд хүний зан үйл, дадал хэвшлийг өөрчлөхөд дижитал технологийг ашиглах асуудал нь зан үйлийн шинжлэх ухаан болон програм хангамжийн инженерийн огтлолцолд байрлах чухал судалгааны чиглэл болж байна. Эрүүл мэнд, боловсрол, бүтээмж, хувийн зохион байгуулалт зэрэг олон талбар нь давтагдах зан үйлд тулгуурладаг бөгөөд ийм тогтворжсон давталтыг дадал гэж үздэг \cite{ref1}. Дадал нь тодорхой нөхцөл, өдөөгч дохио, давталтын нөлөөгөөр аажмаар автоматжиж, хүний өдөр тутмын шийдвэр гаргалт, үйл ажиллагаанд тогтвортой нөлөө үзүүлдэг. Иймээс зан үйлийг тогтвортой өөрчлөхөд зөвхөн богино хугацааны сэдэл, сануулга хангалтгүй бөгөөд дадал үүсэх суурь механизмыг ойлгосон системийн шийдэл шаардлагатай.

Дадал үүсэх үйл явцыг тайлбарлахад Duhigg-ийн өдөөгч дохио--үйлдлийн хэвшил--урамшууллын загвар, Clear-ийн өдөөгч дохио--хүсэл тэмүүлэл--хариу үйлдэл--урамшууллын загвар, мөн Fogg-ийн зан үйлийн загвар өргөн хэрэглэгддэг \cite{ref2,ref3,ref4}. Эдгээр онол нь зан үйлийг ямар нөхцөлд эхлүүлэх, тогтвортой давтуулах, улмаар автоматжуулахыг тайлбарладаг тул дадал төлөвшүүлэх дижитал системийн логикийг тодорхойлоход онолын суурь болдог.

Өнөөгийн дадал хянах олон систем сануулга өгөх, гүйцэтгэл бүртгэх, цуваа болон статистик мэдээлэл харуулах зэрэг нийтлэг функцтэй боловч дадал үүсэх зан үйлийн механизмыг цогц байдлаар тусгаж чаддаггүй. Судалгаагаар ийм төрлийн шийдлүүд нь хэрэглэгчийг давталтад дэмжих боловч зарим тохиолдолд системийн сануулга, интерфэйсээс хэт хамааралтай болгох эрсдэлтэйг тэмдэглэсэн байдаг \cite{ref5}. Үүнээс гадна зан үйлийн онолыг системийн архитектур, өгөгдлийн загвар, функциональ логик, алгоритмын түвшинд хөрвүүлэн хэрэгжүүлэх инженерчлэлийн судалгаа харьцангуй хязгаарлагдмал байна.

Иймд энэхүү судалгаагаар дадал төлөвшлийн мөчлөгт тулгуурлан хэрэглэгчийн дадал үүсэх, бэхжих үйл явцыг дэмжих зан үйлд суурилсан програм хангамжийн системийн загвар боловсруулна. Уг системд өдөөгч дохио, хүсэл тэмүүлэл, хариу үйлдэл, урамшуулал гэсэн бүрэлдэхүүнүүдийг системийн функц, өгөгдлийн загвар, алгоритмын логиктой уялдуулан авч үзнэ. Өөрөөр хэлбэл, энэ ажил нь зөвхөн дадал бүртгэх энгийн систем бус, харин хэрэглэгчийн зан үйлийг нөхцөл байдалд суурилсан сануулга, сэдлийн дэмжлэг, төвөгшил багатай гүйцэтгэл, оновчтой урамшууллын механизмаар дамжуулан тогтвортой төлөвшүүлэх системийн шийдлийг санал болгож байгаагаараа онцлог юм.

\section{Судалгааны зорилго}

Энэхүү судалгааны ажлын зорилго нь дадал төлөвшлийн мөчлөгт тулгуурлан хэрэглэгчийн дадал хэвшил тогтох үйл явцыг дэмжих зан үйлд суурилсан програм хангамжийн системийн загвар, архитектур, системийн зохиомж, туршилтын загварыг боловсруулахад оршино.

\section{Судалгааны зорилтууд}

Дээрх зорилгыг хэрэгжүүлэхийн тулд дараах зорилтуудыг дэвшүүлэв.

\begin{enumerate}
    \item Дадал үүсэх механизм, дадал төлөвшлийн мөчлөг, Фоггийн зан үйлийн загвар, итгүүлэн нөлөөлөх технологийн онолын үндсийг судлах.
    \item Одоо ашиглагдаж буй дадал хянах системүүдийн онцлог, давуу тал, хязгаарлалтыг харьцуулан шинжлэх.
    \item Зан үйлд суурилсан дадал төлөвшүүлэх системийн үзэл баримтлал, функциональ болон функциональ бус шаардлагыг тодорхойлох.
    \item Өдөөгч дохио, хүсэл тэмүүлэл, хариу үйлдэл, урамшуулал зэрэг бүрэлдэхүүнүүдийг системийн функц, өгөгдлийн загвар, алгоритмын логиктой уялдуулан загварчлах.
    \item Дадлын хүчийг автоматжилтын хэмжүүр болон зан үйлийн үзүүлэлтүүдийг хослуулан үнэлэх модель боловсруулах.
    \item Боловсруулсан шийдэлд тулгуурлан системийн туршилтын загварыг хэрэгжүүлж, үр нөлөөг үнэлэх.
\end{enumerate}

\section{Судалгааны асуулт}

Энэхүү судалгааны ажил нь дараах үндсэн асуултуудад хариулахыг зорьж байна.

\textbf{СА-1.} Дадал төлөвшлийн мөчлөг, Фоггийн зан үйлийн загвар, итгүүлэн нөлөөлөх системийн зохиомжийн зарчмуудыг зан үйлд суурилсан програм хангамжийн системийн архитектур, өгөгдлийн загвар, үндсэн алгоритмын түвшинд хэрхэн хэрэгжүүлж болох вэ?

\textbf{СА-2.} Дадлын хүчийг автоматжилтын хэмжүүр болон зан үйлийн дэмжих үзүүлэлтүүдийг хослуулан хэрхэн үнэлж болох вэ?

\textbf{СА-3.} Нөхцөл байдалд суурилсан сануулга, төвөгшлийг бууруулах дэмжлэг, эргэцүүлсэн урамшууллыг хослуулсан туршилтын загвар нь хэрэглэгчийн дадал төлөвшилд хэрхэн нөлөөлөх вэ?

\section{Судалгааны ач холбогдол}

Энэхүү ажил нь онолын, практик болон инженерчлэлийн ач холбогдолтой.

Онолын ач холбогдлын хувьд, дадал үүсэх зан үйлийн онолыг програм хангамжийн системийн загварчлалтай уялдуулан авч үзэж, дадал төлөвшлийн мөчлөг болон дадлын хүчний ойлголтыг системийн логикт буулгах оролдлого хийснээрээ ач холбогдолтой.

Практик ач холбогдлын хувьд, уг судалгааны үр дүнд боловсруулагдах системийн загвар нь эрүүл хэвшил, суралцах дадал, бүтээмж, хувийн зохион байгуулалт зэрэг зан үйлийг дэмжих дижитал шийдэл боловсруулах суурь болно.

Инженерчлэлийн ач холбогдлын хувьд, зан үйлийн онолыг системийн шаардлага, өгөгдлийн бүтэц, алгоритм, архитектур, системийн зохиомжтой уялдуулан хэрэгжүүлэх жишиг загвар болох боломжтой.

\section{Судалгааны арга зүй}

Судалгаанд дараах үндсэн аргуудыг ашиглана.

Нэгдүгээрт, онолын судалгааны аргаар дадал үүсэх онол, зан үйлийн өөрчлөлтийн загвар, итгүүлэн нөлөөлөх технологи, дадлын хүчийг хэмжих аргуудтай холбоотой бүтээлүүдийг судалж, онолын үндэслэлийг тодорхойлно.

Хоёрдугаарт, харьцуулсан шинжилгээний аргаар одоо ашиглагдаж буй дадал хянах системүүдийн онцлог, давуу тал, хязгаарлалтыг шинжилнэ.

Гуравдугаарт, системийн загварчлалын аргаар онолын хандлагуудыг системийн архитектур, өгөгдлийн бүтэц, функциональ шаардлага, алгоритмын логиктой уялдуулан загварчилна.

Дөрөвдүгээрт, програм хангамжийн хөгжүүлэлтийн аргаар боловсруулсан шаардлага, архитектур, системийн зохиомжийн дагуу туршилтын загварыг хэрэгжүүлнэ.

Тавдугаарт, туршилт ба үнэлгээний аргаар системийн хэрэглээний өгөгдөл, гүйцэтгэлийн бүртгэл, автоматжилтын үнэлгээ, хэрэглэгчийн оролцоонд тулгуурлан системийн үр нөлөөг үнэлнэ.